\section{Conclusioni}
In questo lavoro è stato studiato il modello di Ising bidimensionale attraverso tre approcci distinti, mostrando 
come i tre approcci si rivelino complementari e sinergici nello studio del modello: 
la teoria di campo medio offre intuizione analitica, le simulazioni MC permettono di quantificare con precisione 
le osservabili termodinamiche, 
e la PCA dimostra il potenziale degli strumenti di data science nell'esplorazione di sistemi fisici complessi.

La teoria di campo medio ha fornito una descrizione analitica della transizione, 
stimando una temperatura critica $T_C^{MF}= \frac{4J}{k_b}$, che sovrastima il valore esatto di Onsager$ T_C \approx 2.269 \frac{J}{k_b}$.
Tale discrepanza è dovuta al fatto che sono state trascurate le fluttuazioni locali degli spin, che risultano particolarmente rilevanti in prossimità della transizione di fase.

Le simulazioni Monte Carlo con algoritmo di Metropolis hanno consentito 
di riprodurre numericamente le principali grandezze termodinamiche del sistema. I risultati ottenuti
sono in buon accordo con la soluzione esatta: le stime della temperatura critica, ricavate 
dal picco della suscettività e del calore specifico sono compatibili con il valore esatto calcolato da Onsager entro 
gli errori stimati.

Infine, l'applicazione della PCA alle configurazioni generate con 
Monte Carlo ha mostrato come tecniche di apprendimento 
automatico non supervisionato siano in grado di identificare 
automaticamente il parametro d'ordine del sistema e la transizione di fase, 
senza alcuna conoscenza a priori della fisica sottostante. La prima componente principale
risulta correlata quasi perfettamente con la magnetizzazione ($r \approx \ 1$), mentre 
la seconda componente riflette il comportamento della suscettività ($ r \approx \ 0.7).
